\scanpage{58}
\mspage{3}
% @ top left, before line 1 | display
\[
 \deleted{X_{//a}\cap X_{\unsure{/b|//b}}}
\]

% @ top, lines 1–4 | text
Soit $X$ loc irréductible [i.e telle que les comp. \deleted{irréductibles} \inserted{connexes} soient irréductibles i.e. satisfassent \deleted{\unsure{pg|$P_S$}} $*1'$]. Alors la relation \unsure{a|:}

% @ upper third, line 5 | display
\[
 R(a,b):\qquad X_{//a}\cap X_{//b}\ne\varnothing
\]

% @ upper third, lines 6–8 | text
est une relation d'équivalence dans $\operatorname{Ob}X$, laquelle n'est autre que celle \deleted{\unsure{les}} définie par

% @ upper middle, line 8 | display
\[
 \operatorname{Ob}X\longrightarrow\pi_0(X).
\]

% @ middle, lines 9–10 and left | text
De plus, si $X_{//a}\cap X_{//b}\deleted{\ne\varnothing},$ est non vide, il est $0$-conn., \deleted{\unsure{donc} que les \unsure{composantes}} \inserted{et de \unsure{$\overline{\mathrm m}$} pour l'intersection d'un nb fini quelconque $\unsure{\geq|>}\unsure{1|2}$ d'ouverts \inserted{de la forme $X_{//a_i}$ dans une $\overline{\mathrm m}$ comp. conn. de $X$.}}
% slanted

% @ middle, lines 11–13 | text
Si $X$ est $0$-connexe, \deleted{alors} pour que $X$ soit $1$-connexe, il suffit que les $X_{//a}$ le soient.
% archive sheet number: 57
